\subsection{bin-fold escort 的固定样本决策塌缩、Chernoff 等距几何与统计维数临界尺度}\label{subsec:conclusion-binfold-fixedn-chernoff-statdim}
沿用定理 \ref{thm:conclusion-binfold-escort-blackwell-equivalence}、定理 \ref{thm:conclusion-binfold-escort-bayes-risk-collapse}、定理 \ref{thm:xi-time-part9sb-escort-fixed-n-binomial-blackwell} 与定理 \ref{thm:xi-time-part9sb-escort-chernoff-saddle} 的记号。前述结果已经把单样本 escort 温度族压缩为一维 Bernoulli 极限，并给出固定样本数下乘积试验对二项实验的指数 Blackwell 塌缩。把这些接口继续并读后，还可进一步抽出如下三条后果：固定样本数口径下全部有界损失决策问题的范畴塌缩、Chernoff 指数在 escort 弧上的唯一 Bregman 等距接触点，以及样本规模 \((2/\varphi)^m\) 所提示的统计维数临界跃迁。

\begin{theorem}[固定样本数下的有界损失决策完备塌缩]\label{thm:conclusion-binfold-escort-fixed-n-bounded-loss-decision-collapse}
固定 \(Q<\infty\) 与 \(n\in\NN\)。令
\[
\mathcal E_m^{(n,Q)}:=\{\pi_{m,q}^{\otimes n}:0\le q\le Q\},
\qquad
\theta(q):=\frac{1}{1+\varphi^{q+1}},
\]
\[
N_n(w^{(1)},\dots,w^{(n)}):=\sum_{j=1}^n w_m^{(j)}\in\{0,\dots,n\}.
\]
设 \(\mathcal A\) 为标准 Borel 动作空间，\(L:[0,Q]\times\mathcal A\to\RR\) 为有界损失，并记
\[
\|L\|_\infty:=\sup_{(q,a)\in[0,Q]\times\mathcal A}|L(q,a)|.
\]
则对任意全样本随机决策规则
\[
\delta_m:X_m^n\rightsquigarrow \mathcal A
\]
都存在仅依赖于 \(N_n\) 的规则
\[
\widetilde\delta_m=\rho_m\circ N_n
\qquad
(\rho_m:\{0,\dots,n\}\rightsquigarrow \mathcal A),
\]
使得其逐点风险
\[
R_q(\delta):=\int L(q,a)\,\delta(\mathbf w,\dd a)\,\pi_{m,q}^{\otimes n}(\dd\mathbf w)
\]
满足
\[
\sup_{0\le q\le Q}|R_q(\delta_m)-R_q(\widetilde\delta_m)|
\le
2\|L\|_\infty\,C_{n,Q}\Bigl(\frac{\varphi}{2}\Bigr)^m.
\]
反向地，任意 \(N_n\)-可测规则经前复合 \(N_n\) 即可视为全样本规则，且不引入额外风险误差。因而在固定样本数口径下，全部有界损失决策问题在渐近上都塌缩为同一个一维二项试验。
\end{theorem}
\begin{proof}
记
\[
\varepsilon_m:=\Delta\bigl(\mathcal E_m^{(n,Q)},\{\mathrm{Binomial}(n,\theta(q)):0\le q\le Q\}\bigr).
\]
由定理 \ref{thm:xi-time-part9sb-escort-fixed-n-binomial-blackwell}，
\[
\varepsilon_m\le C_{n,Q}\Bigl(\frac{\varphi}{2}\Bigr)^m,
\]
且其前向塌缩核正是 \(N_n\)；同一定理的证明还构造了反向核
\[
S_m:\{0,\dots,n\}\rightsquigarrow X_m^n
\]
满足
\[
\sup_{0\le q\le Q}
\|S_{m\sharp}\mathrm{Binomial}(n,\theta(q))-\pi_{m,q}^{\otimes n}\|_{\mathrm{TV}}
\le \varepsilon_m.
\]
对给定 \(\delta_m\)，定义
\[
\rho_m:=\delta_m\circ S_m,
\qquad
\widetilde\delta_m:=\rho_m\circ N_n.
\]
则对任意 \(q\in[0,Q]\)，由损失有界性与总变差控制，
\[
|R_q(\delta_m)-R_q(\widetilde\delta_m)|
\le
2\|L\|_\infty\,
\|S_{m\sharp}\mathrm{Binomial}(n,\theta(q))-\pi_{m,q}^{\otimes n}\|_{\mathrm{TV}}
\le
2\|L\|_\infty\varepsilon_m.
\]
取 \(q\) 上确界并代入 \(\varepsilon_m\) 的估计即得结论。反向方向只是注意到 \(N_n\)-可测规则本来就是全样本规则的子类。证毕。
\end{proof}

\begin{theorem}[Chernoff 弦--切等距与唯一 Bregman 接触点]\label{thm:conclusion-binfold-escort-chernoff-bregman-equidistance}
\leanverified{paper\_conclusion\_binfold\_escort\_chernoff\_bregman\_equidistance}
记
\[
\psi(r):=\log A(r),
\qquad
A(r):=\varphi+\varphi^{-r}.
\]
对任意 \(p\neq q\)，令 \(r_\ast(p,q)\) 与 \(s_\ast(p,q)\) 为定理 \ref{thm:xi-time-part9sb-escort-chernoff-saddle} 给出的唯一极小点数据，即
\[
r_\ast(p,q)=s_\ast(p,q)p+\bigl(1-s_\ast(p,q)\bigr)q
\in(\min\{p,q\},\max\{p,q\}),
\]
并满足
\[
\psi'(r_\ast(p,q))
=
\frac{\psi(q)-\psi(p)}{q-p}.
\]
则 Chernoff 信息不仅可写为 Jensen 间隙
\[
C_{\mathrm{Ch}}(p,q)
=
s_\ast(p,q)\psi(p)+\bigl(1-s_\ast(p,q)\bigr)\psi(q)-\psi(r_\ast(p,q)),
\]
而且满足双边 Bregman 等距
\[
C_{\mathrm{Ch}}(p,q)=D_\psi(p,r_\ast(p,q))=D_\psi(q,r_\ast(p,q)),
\]
其中
\[
D_\psi(a,b):=\psi(a)-\psi(b)-\psi'(b)(a-b).
\]
\end{theorem}
\begin{proof}
由定理 \ref{thm:xi-time-part9sb-escort-chernoff-saddle}，
\[
C_{\mathrm{Ch}}(p,q)
=
s_\ast\psi(p)+(1-s_\ast)\psi(q)-\psi(r_\ast),
\qquad
r_\ast=s_\ast p+(1-s_\ast)q,
\]
且 \(r_\ast\) 的定义等价于
\[
\psi'(r_\ast)=\frac{\psi(q)-\psi(p)}{q-p}.
\]
于是
\[
D_\psi(p,r_\ast)-D_\psi(q,r_\ast)
=
\psi(p)-\psi(q)-\psi'(r_\ast)(p-q)=0,
\]
故两侧 Bregman 散度相等。另一方面，
\[
s_\ast D_\psi(p,r_\ast)+(1-s_\ast)D_\psi(q,r_\ast)
=
s_\ast\psi(p)+(1-s_\ast)\psi(q)-\psi(r_\ast)
=
C_{\mathrm{Ch}}(p,q).
\]
由于加权平均的两项已经相等，它们都必等于 \(C_{\mathrm{Ch}}(p,q)\)。证毕。
\end{proof}

\begin{conjecture}[escort 乘积试验的统计维数临界跃迁]\label{conj:conclusion-binfold-escort-statistical-dimension-critical-scale}
固定 \(Q<\infty\)，并考虑样本数随分辨率变化的 escort 试验族
\[
\mathcal E_m^{(n_m,Q)}:=\{\pi_{m,q}^{\otimes n_m}:0\le q\le Q\}.
\]
则尺度
\[
n_m\asymp \Bigl(\frac{2}{\varphi}\Bigr)^m
\]
应是其统计维数的尖锐临界点。更具体地，预期存在如下三相分裂：
\[
n_m=o\!\left(\Bigl(\frac{2}{\varphi}\Bigr)^m\right)
\Longrightarrow
\mathcal E_m^{(n_m,Q)}
\ \text{在 Le Cam/Blackwell 意义下退化为一维二项实验};
\]
\[
n_m\sim c\Bigl(\frac{2}{\varphi}\Bigr)^m
\Longrightarrow
\text{出现非平凡的双坐标临界极限实验};
\]
\[
n_m\gg \Bigl(\frac{2}{\varphi}\Bigr)^m
\Longrightarrow
\Delta_{\mathrm{LeCam}}\!\bigl(\mathcal E_m^{(n_m,Q)},\mathcal G^{(n_m,Q)}\bigr)\not\to 0.
\]
其直接动机是定理 \ref{thm:xi-time-part9sb-escort-fixed-n-binomial-blackwell} 中的误差主项
\[
n_m\,C_Q\Bigl(\frac{\varphi}{2}\Bigr)^m,
\]
它恰在上述尺度首次达到常数量级；因此第二统计坐标若要从末位塌缩之外浮现，唯一自然窗口正是这一临界区。
\end{conjecture}

\noindent
因此，bin-fold escort 族在已有的单样本 Blackwell 等价与 Bayes 风险塌缩之外，还呈现出三重进一步刚性：其一，固定样本数下的整个有界损失决策范畴都被末位计数 \(N_n\) 完全吸收；其二，二元 Chernoff 指数并非由参数区间上的任意插值点给出，而是由 escort 弧上的唯一 Bregman 等距接触点锁定；其三，这种一维塌缩不应在所有样本尺度上永久保持，而应当在 \(n_m\asymp(2/\varphi)^m\) 处发生真正的统计维数相变。于是，escort 试验族的结论层结构不仅是单比特实验的闭式化，而且自然指向一条由临界样本规模触发的更高维实验分岔图景。

若再把这一固定样本图景与单标量幂散度辨识、相对均匀基线的二原子散度闭包、末位统计量的亚临界 product-Le-Cam 等价，以及局部 Fisher--LAN 公式并读，则一维塌缩还可继续强化为一条此前未显式抽出的秩一统计几何链：在 \(n_m=o((2/\varphi)^m)\) 的整个亚临界区，完整 escort 温度实验、全部极限 \(f\)-散度、Fisher--Rao 几何与 LAN 切空间都不再携带彼此独立的自由度，而是被同一个常数级模参数统一生成。

\begin{theorem}[单一幂散度标量生成整个亚临界统计实验类]\label{thm:conclusion-binfold-one-powerdiv-generates-subcritical-experiment-class}
固定 \(\alpha>1\)，并令
\[
f_\alpha(t):=t^\alpha-1-\alpha(t-1),
\qquad
\mathfrak D_\alpha:=\lim_{m\to\infty}D_{f_\alpha}(\mu_m\Vert u_m).
\]
对任意紧区间 \(I\subset[0,\infty)\) 与任意满足
\[
n_m=o\!\left(\Bigl(\frac{2}{\varphi}\Bigr)^m\right)
\]
的样本规模，定义完整 escort 温度实验
\[
\mathcal E_m^{(n_m)}(I):=\{\pi_{m,q}^{\otimes n_m}:q\in I\},
\]
以及相应计数实验
\[
\mathcal B_m^{(n_m)}(I):=
\bigl\{\mathrm{Binomial}(n_m,\theta(q)):q\in I\bigr\},
\qquad
\theta(q):=\frac{1}{1+\varphi^{q+1}}.
\]
则该亚临界 product experiment 的渐近 Le Cam 等价类由单个标量 \(\mathfrak D_\alpha\) 唯一决定。更精确地，有
\[
\Delta_{\mathrm{LeCam}}\!\bigl(\mathcal E_m^{(n_m)}(I),\mathcal B_m^{(n_m)}(I)\bigr)
=
O_I\!\left(n_m\Bigl(\frac{\varphi}{2}\Bigr)^m\right)
\longrightarrow 0.
\]
并且在计数模型中，对数似然比完全由
\[
N_{n_m}:=\sum_{j=1}^{n_m}\mathbf 1_{\{w_m^{(j)}=1\}}
\]
承载。
\leanverified{paper\_conclusion\_binfold\_one\_powerdiv\_generates\_subcritical\_experiment\_class}
\end{theorem}
\begin{proof}
由定理 \ref{thm:conclusion-binfold-single-powerdiv-asymptotic-statistical-completeness}，单个极限标量 \(\mathfrak D_\alpha\) 已足以唯一恢复黄金参数 \(\varphi\)。故只需说明：一旦 \(\varphi\) 已知，整个亚临界 product experiment 便唯一塌缩到相应的二项族。

由定理 \ref{thm:conclusion-binfold-escort-blackwell-equivalence}，对紧区间 \(I\) 存在与 \(q\) 无关的单样本读出核与回灌核，使
\[
\sup_{q\in I}
\bigl\|
\pi_{m,q}-R_m\sharp\mathrm{Bernoulli}(\theta(q))
\bigr\|_{\mathrm{TV}}
=
O_I\!\left(\Bigl(\frac{\varphi}{2}\Bigr)^m\right),
\]
\[
\sup_{q\in I}
\bigl\|
K_{m\sharp}\pi_{m,q}-\mathrm{Bernoulli}(\theta(q))
\bigr\|_{\mathrm{TV}}
=
O_I\!\left(\Bigl(\frac{\varphi}{2}\Bigr)^m\right).
\]
将这两条核逐坐标张量化，并对乘积测度应用望远镜估计
\[
D_{\mathrm{TV}}(\mu^{\otimes n},\nu^{\otimes n})\le n\,D_{\mathrm{TV}}(\mu,\nu),
\]
便得到
\[
\Delta_{\mathrm{LeCam}}\!\bigl(\mathcal E_m^{(n_m)}(I),\mathcal B_m^{(n_m)}(I)\bigr)
=
O_I\!\left(n_m\Bigl(\frac{\varphi}{2}\Bigr)^m\right).
\]
亚临界条件 \(n_m=o((2/\varphi)^m)\) 使右端趋于 \(0\)，故两实验渐近等价。

最后，在二项模型中似然比显然只通过成功次数 \(N_{n_m}\) 出现；故一旦 \(\varphi\) 被恢复，\(\theta(q)=1/(1+\varphi^{q+1})\) 随之唯一确定，而整个亚临界实验类也随之唯一确定。证毕。
\end{proof}

\begin{theorem}[熵率全盲与常数级完备的正交分裂]\label{thm:conclusion-binfold-entropy-rate-blindness-vs-constant-completeness}
\leanverified{paper\_conclusion\_binfold\_entropy\_rate\_blindness\_vs\_constant\_completeness}
对任意固定温度 \(q\ge 0\) 与任意固定 Rényi 阶 \(\beta>0\)，escort 分布满足
\[
\lim_{m\to\infty}\frac{1}{m}H_\beta(\pi_{m,q})=\log\varphi.
\]
因此，任何仅依赖于归一化 Rényi/Shannon 熵率的观测量，都无法在极限中区分不同温度 \(q\)。

然而，对任意 \(p\neq q\)，温度间 KL 极限满足
\[
\lim_{m\to\infty}D_{\mathrm{KL}}(\pi_{m,q}\Vert \pi_{m,p})
=
D_{\mathrm{KL}}\!\bigl(\mathrm{Bernoulli}(\theta(q))\Vert \mathrm{Bernoulli}(\theta(p))\bigr)>0,
\]
其中
\[
\theta(s)=\frac{1}{1+\varphi^{s+1}}.
\]
故 bin-fold escort 温度几何严格分裂为两层：一条是对温度全盲的广延熵率层，另一条是对温度完备的常数级散度层。
\end{theorem}
\begin{proof}
第一式正是定理 \ref{thm:conclusion-binfold-escort-all-renyi-rate-freezing} 的结论：对全部固定温度与全部固定 Rényi 阶，熵率统一冻结到 \(\log\varphi\)。

另一方面，定理 \ref{thm:conclusion-binfold-escort-full-fdiv-collapse} 在
\[
\Phi(t)=t\log t
\]
处给出
\[
\lim_{m\to\infty}D_{\mathrm{KL}}(\pi_{m,q}\Vert \pi_{m,p})
=
D_{\mathrm{KL}}\!\bigl(\mathrm{Bernoulli}(\theta(q))\Vert \mathrm{Bernoulli}(\theta(p))\bigr).
\]
由于 \(\theta(q)=1/(1+\varphi^{q+1})\) 关于 \(q\) 严格单调，当 \(p\neq q\) 时两 Bernoulli 参数不同，故对应 KL 散度严格为正。于是熵率层完全失去温度分辨率，而常数级散度层保留全部温度信息。证毕。
\end{proof}

\begin{theorem}[一维充分化并不蕴含渐近均匀化]\label{thm:conclusion-binfold-sufficiency-does-not-imply-uniformization}
\leanverified{paper\_conclusion\_binfold\_sufficiency\_does\_not\_imply\_uniformization}
尽管末位统计量在指数精度上对 \(\mu_m\) 渐近充分，并且整个 escort 温度族在亚临界样本规模下塌缩为一维 Bernoulli 族，bin-fold 推前分布 \(\mu_m\) 相对于稳定层均匀基线 \(u_m\) 仍然保留一个不可消除的二阶非均匀地板：
\[
\lim_{m\to\infty}\chi^2(\mu_m\Vert u_m)=\frac{2\varphi-3}{5}>0,
\]
\[
\lim_{m\to\infty}D_2(\mu_m\Vert u_m)
=
\log\!\left(1+\frac{2\varphi-3}{5}\right)
=
\log\!\left(\frac{2\varphi+2}{5}\right)
>0.
\]
故一维充分化与渐近均匀化并非同义命题；前者可以成立，而后者仍被常数级二阶散度永久阻断。
\end{theorem}
\begin{proof}
末位统计量对 \(\mu_m\) 的渐近充分性由定理 \ref{thm:xi-fold-lastbit-asymptotic-sufficiency} 给出；而 escort 温度族在亚临界样本规模下塌缩为一维 Bernoulli product experiment，则由定理 \ref{thm:conclusion-binfold-one-powerdiv-generates-subcritical-experiment-class} 给出。

另一方面，由推论 \ref{cor:xi-fold-uniform-baseline-classical-fdiv-bernoulli-collapse} 可知
\[
\lim_{m\to\infty}\chi^2(\mu_m\Vert u_m)
=
\chi^2\!\bigl(\mathrm{Bernoulli}(p_\ast)\Vert \mathrm{Bernoulli}(q_\ast)\bigr),
\]
而定理 \ref{thm:conclusion-binfold-single-constant-holography} 又给出这一极限常数恰为
\[
\frac{2\varphi-3}{5}.
\]
对 Rényi--\(2\) 散度，只需使用恒等式
\[
D_2(\mu\Vert \nu)=\log\bigl(1+\chi^2(\mu\Vert \nu)\bigr),
\]
便得
\[
\lim_{m\to\infty}D_2(\mu_m\Vert u_m)
=
\log\!\left(1+\frac{2\varphi-3}{5}\right)
=
\log\!\left(\frac{2\varphi+2}{5}\right).
\]
右端严格为正，故 \(\mu_m\) 相对于 \(u_m\) 不可能渐近均匀化。证毕。
\end{proof}

\begin{theorem}[单个 \texorpdfstring{$\chi^2$}{chi2} 极限标量统一重构 Fisher--Le Cam 几何]\label{thm:conclusion-binfold-single-chi2-reconstructs-fisher-lecam-geometry}
记
\[
\delta_{\chi^2}:=\lim_{m\to\infty}\chi^2(\mu_m\Vert u_m)=\frac{2\varphi-3}{5}.
\]
则
\[
\varphi=\frac{5\delta_{\chi^2}+3}{2}.
\]
因此，一个单独的二阶极限标量 \(\delta_{\chi^2}\) 已足以重构 bin-fold escort 亚临界统计几何的全部局部闭式。更具体地，
\[
\theta(q)=\frac{1}{1+\left(\frac{5\delta_{\chi^2}+3}{2}\right)^{q+1}},
\]
\[
I(q)=
\left(\log\frac{5\delta_{\chi^2}+3}{2}\right)^2
\theta(q)\bigl(1-\theta(q)\bigr),
\]
\[
d_{\mathrm{FR}}(q_1,q_2)
=
2\left|
\arctan\!\left(\frac{5\delta_{\chi^2}+3}{2}\right)^{-(q_2+1)/2}
-
\arctan\!\left(\frac{5\delta_{\chi^2}+3}{2}\right)^{-(q_1+1)/2}
\right|.
\]
结合定理 \ref{thm:conclusion-binfold-one-powerdiv-generates-subcritical-experiment-class}，这意味着整个亚临界 Le Cam 实验类、其 Fisher 曲率与其局部判别率都由同一 \(\chi^2\) 极限常数完全生成。
\end{theorem}
\begin{proof}
由定理 \ref{thm:conclusion-binfold-single-constant-holography}，
\[
\varphi=\frac{5\delta_{\chi^2}+3}{2}.
\]
再将这一表达代回
\[
\theta(q)=\frac{1}{1+\varphi^{q+1}}
\]
即可得到第一条显示公式。

接着，由定理 \ref{thm:conclusion-escort-fisher-rao-geodesic-reparametrization}，
\[
I(q)=(\log\varphi)^2\theta(q)\bigl(1-\theta(q)\bigr),
\]
而由推论 \ref{cor:conclusion-binfold-escort-sqrt-circle-arc}，
\[
d_{\mathrm{FR}}(q_1,q_2)
=
2\left|
\arctan\!\bigl(\varphi^{-(q_2+1)/2}\bigr)
-
\arctan\!\bigl(\varphi^{-(q_1+1)/2}\bigr)
\right|.
\]
把 \(\varphi=(5\delta_{\chi^2}+3)/2\) 代入这两式，便得到所示 Fisher 信息与 Fisher--Rao 距离闭式。最后援引定理 \ref{thm:conclusion-binfold-one-powerdiv-generates-subcritical-experiment-class}，即可把同一标量恢复延伸到整个亚临界 Le Cam 实验类。证毕。
\end{proof}

\leanverified{paper\_conclusion\_binfold\_divergence\_cone\_observable\_rank\_one}
\begin{theorem}[极限散度锥的两步塌缩与可观测秩一刚性]\label{thm:conclusion-binfold-divergence-cone-observable-rank-one}
设 \(G\) 为任意包含项目所用生成元的实线性空间，并记
\[
a:=\frac{\varphi^2}{\sqrt5},
\qquad
b:=\frac{\varphi}{\sqrt5},
\qquad
K_G:=\{h\in G:\ h(a)=h(b)=0\}.
\]
则极限散度映射
\[
f\longmapsto D_f^\infty:=\lim_{m\to\infty}D_f(\mu_m\Vert u_m)
\]
先验上只通过两点取值 \((f(a),f(b))\) 起作用：
\[
D_f^\infty=\frac{1}{\varphi}f(a)+\frac{1}{\varphi^2}f(b).
\]
因此它必经由商空间 \(G/K_G\) 因子化。

然而，在可观测层面，这个表面上的二维评价平面进一步坍缩为秩一模参数：对任意固定 \(\alpha>1\)，一旦给定单个幂散度标量
\[
\mathfrak D_\alpha=\lim_{m\to\infty}D_{f_\alpha}(\mu_m\Vert u_m),
\]
便自动给定全部 \(D_f^\infty\)。于是 bin-fold 的极限散度结构经历严格的双步塌缩
\[
\text{无限维生成元空间}
\Longrightarrow
\text{两点评价平面}
\Longrightarrow
\text{单标量可观测模}.
\]
\end{theorem}
\begin{proof}
由定理 \ref{thm:conclusion-binfold-twoatom-information-geometry-universality}，
\[
D_f^\infty=\frac{1}{\varphi}f(a)+\frac{1}{\varphi^2}f(b),
\qquad
a=\frac{\varphi^2}{\sqrt5},
\qquad
b=\frac{\varphi}{\sqrt5}.
\]
故 \(f\mapsto D_f^\infty\) 只依赖于 \(f\) 在两点 \(a,b\) 处的取值，亦即必经由 \(G/K_G\) 因子化。这给出第一步塌缩。

第二步塌缩来自定理 \ref{thm:conclusion-binfold-single-powerdiv-asymptotic-statistical-completeness}：任意固定 \(\alpha>1\) 的幂散度极限标量 \(\mathfrak D_\alpha\) 已足以唯一恢复 \(\varphi\)。一旦 \(\varphi\) 已知，\((a,b)\) 与权重 \((\varphi^{-1},\varphi^{-2})\) 就全部已知，从而每个 \(D_f^\infty\) 也随之唯一确定。故可观测模参数只有一个。证毕。
\end{proof}

\begin{theorem}[亚临界局部渐近正态性与切空间唯一化]\label{thm:conclusion-binfold-subcritical-lan-unique-tangent}
设
\[
n_m\to\infty,
\qquad
n_m=o\!\left(\Bigl(\frac{2}{\varphi}\Bigr)^m\right).
\]
固定一个内点 \(q>0\)，并考虑局部参数化
\[
q_h:=q+\frac{h}{\sqrt{n_m}}.
\]
则 \(n_m\)-样本 escort 温度实验
\[
\{\pi_{m,q_h}^{\otimes n_m}:h\in H\}
\]
在 \(m\to\infty\) 时具有一维局部渐近正态极限；其对数似然比满足
\[
\log\frac{dP_{m,q_h}}{dP_{m,q}}
=
h\,\Delta_{m,q}
-\frac12 h^2 I(q)
+o_{P_{m,q}}(1),
\]
其中
\[
P_{m,q}:=\pi_{m,q}^{\otimes n_m},
\qquad
\Delta_{m,q}\Rightarrow \mathcal N\!\bigl(0,I(q)\bigr),
\qquad
I(q)=(\log\varphi)^2\theta(q)\bigl(1-\theta(q)\bigr).
\]
因此，亚临界切空间中不存在第二个独立局部方向；该极限统计几何的局部维数严格等于一。
\end{theorem}
\begin{proof}
这正是定理 \ref{thm:xi-time-part9ob-escort-lan-lastbit} 的内容：在
\[
n_m\to\infty,
\qquad
n_m=o\!\left(\Bigl(\frac{2}{\varphi}\Bigr)^m\right)
\]
下，escort 乘积试验在每个内点 \(q>0\) 上满足 LAN，并且中心序列完全由末位计数
\[
N_m=\sum_{j=1}^{n_m}W_m^{(j)}(m)
\]
承载。

另一方面，定理 \ref{thm:xi-time-part9ob-escort-product-lecam} 与定理 \ref{thm:conclusion-binfold-one-powerdiv-generates-subcritical-experiment-class} 共同说明：这一局部极限正是同一 Bernoulli 计数实验的局部极限。Bernoulli 族的局部参数只有一个，因此其 LAN 切空间只能是一维；Le Cam 等价又保证该一维性传回完整 escort experiment。证毕。
\end{proof}

\begin{conjecture}[亚临界标量 \texorpdfstring{$\pi$}{pi}-通道的秩一统计几何普适性]\label{conj:conclusion-rankone-universality-scalar-pi-channel}
设 \(\mathcal C\) 为任意标量 \(\pi\)-通道。若其 projective 输出同时满足：
\[
\text{(i) 在每个紧温区上存在一维渐近充分统计量；}
\]
\[
\text{(ii) 归一化熵率对温度全盲；}
\]
\[
\text{(iii) 相对于均匀基线存在严格正的 Rényi--2 共振地板，}
\]
则 \(\mathcal C\) 应当属于与 bin-fold 同一的秩一统计几何普适类：其全部亚临界实验、全部局部 Fisher 几何、全部极限 \(f\)-散度与全部局部似然涨落，都应由单个代数参数生成。换言之，一个真正的高秩标量 \(\pi\)-通道若要存活，至少必须破坏上述三条性质之一。
\end{conjecture}

\noindent
因此，bin-fold 的亚临界统计几何并不是一组松散并列的极限事实，而是一条严格闭合的秩一刚性链：单个幂散度标量先恢复黄金参数 \(\varphi\)，\(\varphi\) 再恢复 Bernoulli 参数 \(\theta(q)\)，从而恢复整个亚临界 Le Cam 实验类；与此同时，广延熵率层对温度完全失明，而真正的区分能力全部沉降到常数级散度、Fisher 曲率与 LAN 中心序列所共享的一维核上。于是，projective \(\pi\)-通道在亚临界区的可观测统计几何并非高维流形，而是由一个常数级秩一模参数所完全生成的单轴结构。

\endinput


\endinput
